\begin{abstract}
Diabetic Retinopathy (DR) remains a major cause of preventable blindness worldwide, 
impacting roughly 35\% of diabetic patients. Federated learning (FL) enables multiple 
hospitals to grade DR collaboratively without sharing sensitive retinal images. 
However, baseline algorithms like FedAvg and FedProx degrade severely under non-IID 
data distributions caused by varying patient demographics across hospital tiers. To 
address this, we introduce \textbf{OrdinalFed}, which incorporates two components: 
(1)~an \emph{OrdinalLoss} that penalizes grade-skipping errors based on their squared 
clinical distance, and (2)~a \emph{QWK-weighted aggregation} scheme that uses a 
cosine-annealing schedule to shift from sample counts to Quadratic Weighted Kappa 
(QWK) validation scores on a 50-image server set. On the DDR benchmark across three 
heterogeneity settings, OrdinalFed matches baseline peak QWK ($\approx$0.799) under 
heavy non-IID skew ($\alpha=0.1$) while raising the worst-case floor to \textbf{0.6238} 
(compared to 0.5963 for FedProx and 0.4064 for FedAvg), which narrows the best-to-worst 
gap to 0.1751 rather than 0.2024 and 0.3929. In a realistic multi-tier hospital partition, 
OrdinalFed achieves a higher peak QWK of \textbf{0.8325} (versus 0.8312 for FedProx 
and 0.8223 for FedAvg) with no added communication overhead.
\end{abstract}

\begin{IEEEkeywords}
Diabetic Retinopathy, Federated Learning, Non-IID,
Ordinal Loss, QWK-Weighted Aggregation, EfficientNet-B3,
Medical Imaging
\end{IEEEkeywords}